\section{لم ناکایاما و رشته‌های دقیق}

% Source: sources/lectures/16 Nakayama Lemma and Exact sequences .pdf | PDF page: 1 | item: Nakayama lemma statement
\begin{frame}[t]{لم ناکایاما}
    \label{lec16:nakayama}
    \begin{theorem}[لم ناکایاما]
        فرض کنید \(R\) یک حلقه و \(I\) یک ایده‌آل چپ از \(R\) باشد که
        \[
            I\subseteq J(R).
        \]
        اگر \(M\) یک \(R\)-مدول با تولید متناهی باشد و
        \[
            IM=M,
        \]
        آنگاه \(M=0\).
    \end{theorem}
\end{frame}

% Source: sources/lectures/16 Nakayama Lemma and Exact sequences .pdf | PDF page: 1 | item: minimal generating set and expression of m_1 in J(R)M
\begin{frame}[t]{اثبات ناکایاما: انتخاب مولدهای کمینه}
    \label{lec16:nakayama-minimal}
    مجموعهٔ مولد
    \[
        M=\langle m_1,m_2,\ldots,m_k\rangle
    \]
    را با کمترین تعداد ممکن انتخاب می‌کنیم. اگر \(M\ne0\)، آنگاه \(k\ge1\).
    چون
    \[
        M=IM\subseteq J(R)M\subseteq M,
    \]
    داریم \(M=J(R)M\). پس برای \(m_1\in M\)،
    \[
        m_1=r_1x_1+\cdots+r_tx_t,
        \qquad r_i\in J(R),\quad x_i\in M.
    \]
\end{frame}

% Source: sources/lectures/16 Nakayama Lemma and Exact sequences .pdf | PDF page: 1 | item: expansion of x_i in the generators, substitution, and two-sidedness footnote
% AMBIGUITY[L16-Q02]: primary, PDF page 1, substitution ends with r_k alpha_{k1} after t summands; the upper limit t is an explicitly provisional index normalization.
\begin{frame}[t]{اثبات ناکایاما: جایگذاری مولدها}
    \label{lec16:nakayama-coefficients}
    هر \(x_i\) را بر حسب مولدها می‌نویسیم:
    \[
        x_i=\alpha_{i1}m_1+\cdots+\alpha_{ik}m_k,
        \qquad 1\le i\le t.
    \]
    با جایگذاری در رابطهٔ \(m_1\)،
    \[
        m_1=r'_1m_1+\beta_2m_2+\cdots+\beta_km_k,
        \qquad r'_1=\sum_{i=1}^{t}r_i\alpha_{i1}.
    \]
    چون \(J(R)\) ایده‌آلی دوطرفه است،
    \[
        r_i\alpha_{i1}\in J(R)\quad\Longrightarrow\quad r'_1\in J(R).
    \]
    {\footnotesize در جزوه بر دوطرفه بودن \(J(R)\) تأکید شده است.
    کرانِ جمع را موقتاً \(t\) خوانده‌ایم؛ جزوه در جایگذاری \(k\) نوشته است: \lr{L16-Q02}.}
\end{frame}

% Source: sources/lectures/16 Nakayama Lemma and Exact sequences .pdf | PDF page: 1 | item: invertibility of 1-r'_1, deletion of a generator, and contradiction
% AMBIGUITY[L16-Q02]: continuation uses the provisional t/k index normalization from the previous frame.
\begin{frame}[t]{اثبات ناکایاما: حذف یک مولد}
    \label{lec16:nakayama-contradiction}
    رابطهٔ قبل می‌دهد
    \[
        (1-r'_1)m_1=\beta_2m_2+\cdots+\beta_km_k.
    \]
    در جزوه ابتدا \(r'_1\ne1\) مطرح شده است. \alert{چرا؟}
    \par\medskip
    چون \(r'_1\in J(R)\)، عنصر \(1-r'_1\) وارون‌پذیر است.
    \alert{چرا؟}
    بنابراین
    \[
        m_1\in\langle m_2,\ldots,m_k\rangle,
        \qquad M=\langle m_2,\ldots,m_k\rangle.
    \]
    این با کمینه بودن تعداد مولدها تناقض دارد؛ پس \(M=0\).
    \hfill\(\qed\)
    \medskip

    {\footnotesize قرارداد شاخص‌ها در گام جایگذاری موقت است: \lr{L16-Q02}.}
\end{frame}

% Source: sources/lectures/16 Nakayama Lemma and Exact sequences .pdf | PDF page: 2 | item: quotient form of Nakayama and its proof
\begin{frame}[t]{نتیجه‌ای از لم ناکایاما}
    \label{lec16:nakayama-quotient}
    \begin{corollary}
        فرض کنید \(N\) زیرمدولی از \(M\)، و \(M/N\) با تولید متناهی باشد.
        اگر \(I\) ایده‌آلی چپ از \(R\) باشد و
        \[
            I\subseteq J(R),\qquad N+IM=M,
        \]
        آنگاه \(M=N\).
    \end{corollary}
    \textbf{اثبات.} روی مدول خارج‌قسمت داریم
    \[
        I(M/N)=\frac{IM+N}{N}=\frac{M}{N}.
    \]
    از لم ناکایاما، \(M/N=0\)؛ بنابراین \(M=N\). \hfill\(\qed\)
\end{frame}

% Source: sources/lectures/16 Nakayama Lemma and Exact sequences .pdf | PDF page: 3 | item: definition of exactness at the middle module
\begin{frame}[t]{رشتهٔ دقیق}
    \label{lec16:exact-at-middle}
    \begin{definition}
        فرض کنید \(A,B,C\) سه \(R\)-مدول و \(f,g\) دو همریختی باشند:
        \[
            A\xrightarrow{f}B\xrightarrow{g}C.
        \]
        رشتهٔ بالا را \textbf{دقیق} می‌گوییم هرگاه
        \[
            \operatorname{Im}f=\ker g.
        \]
    \end{definition}
\end{frame}

% Source: sources/lectures/16 Nakayama Lemma and Exact sequences .pdf | PDF page: 3 | item: injectivity and surjectivity as exactness
\begin{frame}[t]{دقت و یک‌به‌یکی یا پوشا بودن}
    \label{lec16:exact-mono-epi}
    \begin{remark}
        اگر رشتهٔ
        \[
            0\xrightarrow{\varphi}A\xrightarrow{f}B
        \]
        دقیق باشد، آنگاه
        \[
            \ker f=\operatorname{Im}\varphi=0;
        \]
        پس \(f\) یک‌به‌یک است.
    \end{remark}
    \begin{remark}
        اگر رشتهٔ \(B\xrightarrow{g}C\xrightarrow{\psi}0\) دقیق باشد،
        \[
            \operatorname{Im}g=\ker\psi=C;
        \]
        پس \(g\) پوشاست.
    \end{remark}
\end{frame}

% Source: sources/lectures/16 Nakayama Lemma and Exact sequences .pdf | PDF page: 3 | item: exact sequence and short exact sequence definitions
\begin{frame}[t]{رشتهٔ دقیق و رشتهٔ دقیق کوتاه}
    \label{lec16:short-exact-definition}
    \begin{definition}
        دنبالهٔ
        \[
            M_1\xrightarrow{f_1}M_2\xrightarrow{f_2}M_3
            \longrightarrow\cdots\xrightarrow{f_{k-1}}M_k
        \]
        را دقیق می‌گوییم هرگاه در هر جایگاه میانی،
        \[
            \operatorname{Im}f_i=\ker f_{i+1}.
        \]
    \end{definition}
    \begin{definition}
        رشتهٔ دقیقِ به صورت
        \[
            0\longrightarrow A\longrightarrow B\longrightarrow C\longrightarrow0
        \]
        را \textbf{رشتهٔ دقیق کوتاه} می‌نامیم.
    \end{definition}
\end{frame}

% Source: sources/lectures/16 Nakayama Lemma and Exact sequences .pdf | PDF page: 3 | item: nZ and submodule quotient short exact sequence examples
\begin{frame}[t]{دو مثال از رشته‌های دقیق کوتاه}
    \label{lec16:short-exact-examples}
    \begin{example}
        با نگاشت شمول و نگاشت طبیعی، رشتهٔ زیر دقیق کوتاه است:
        \[
            0\longrightarrow n\mathbb Z\lhook\joinrel\longrightarrow
            \mathbb Z\longrightarrow\mathbb Z_n\longrightarrow0.
        \]
    \end{example}
    \begin{example}
        اگر \(N\) زیرمدولی از \(R\)-مدولِ \(M\) باشد، آنگاه
        \[
            0\longrightarrow N\lhook\joinrel\longrightarrow M
            \longrightarrow M/N\longrightarrow0
        \]
        یک رشتهٔ دقیق کوتاه است.
    \end{example}
\end{frame}

% Source: sources/lectures/16 Nakayama Lemma and Exact sequences .pdf | PDF page: 4 | item: short exact sequences attached to a monomorphism and an epimorphism
\begin{frame}[t]{مثال‌های حاصل از یک همریختی}
    \label{lec16:mono-epi-examples}
    \begin{example}
        اگر \(f:M\to N\) یک مونومورفیسم، یعنی همریختی یک‌به‌یک، باشد،
        \[
            0\longrightarrow M\xrightarrow{f}N
            \longrightarrow N/\operatorname{Im}f\longrightarrow0
        \]
        دقیق است؛ نگاشت بعد از \(f\)، نگاشت طبیعی است.
    \end{example}
    \begin{example}
        اگر \(f:M\to N\) یک اپی‌مورفیسم، یعنی همریختی پوشا، باشد،
        \[
            0\longrightarrow\ker f\longrightarrow M\xrightarrow{f}N
            \longrightarrow0
        \]
        دقیق است.
    \end{example}
\end{frame}

% Source: sources/lectures/16 Nakayama Lemma and Exact sequences .pdf | PDF page: 4 | item: short five lemma diagram and three assertions
\begin{frame}[t]{لم پنج‌تایی کوتاه}
    \label{lec16:short-five-lemma}
    فرض کنید نمودار زیر از \(R\)-مدول‌ها و همریختی‌ها جابجایی باشد و هر دو سطر آن دقیق کوتاه باشند:
    \[
    \begin{tikzcd}[ampersand replacement=\&,column sep=large,row sep=large]
        0\arrow[r]\& M\arrow[r,"f"]\arrow[d,"\alpha"']
        \& N\arrow[r,"g"]\arrow[d,"\beta"]
        \& K\arrow[r]\arrow[d,"\gamma"]\&0\\
        0\arrow[r]\& M'\arrow[r,"\varphi"']
        \& N'\arrow[r,"\psi"']\& K'\arrow[r]\&0
    \end{tikzcd}
    \]
    \begin{enumerate}
        \item اگر \(\alpha,\gamma\) یک‌به‌یک باشند، \(\beta\) یک‌به‌یک است.
        \item اگر \(\alpha,\gamma\) پوشا باشند، \(\beta\) پوشاست.
        \item اگر \(\alpha,\gamma\) یکریختی باشند، \(\beta\) یکریختی است.
    \end{enumerate}
\end{frame}

% Source: sources/lectures/16 Nakayama Lemma and Exact sequences .pdf | PDF pages: 4-5 | item: proof of injectivity, from beta(n)=0 to n in Im f; red map-name inconsistency
% AMBIGUITY[L16-Q01]: primary, PDF page 5, red proof after n in ker g = Im f writes g(m)=n; f(m)=n is used provisionally, not silently corrected.
\begin{frame}[t]{لم پنج‌تایی کوتاه: یک‌به‌یکیِ \(\beta\)}
    \label{lec16:five-injective-preimage}
    فرض کنید \(\alpha,\gamma\) یک‌به‌یک باشند. برای \(n\in N\) با \(\beta(n)=0\)،
    \[
        0=\psi(0)=\psi\beta(n)=\gamma g(n).
    \]
    چون \(\ker\gamma=0\)، داریم
    \[
        g(n)=0\quad\Longrightarrow\quad n\in\ker g=\operatorname{Im}f.
    \]
    \begin{alertblock}{خوانش موقتِ یک نماد: \lr{L16-Q01}}
        در این گام، رابطهٔ انتخاب پیش‌تصویر را به صورت
        \(\exists m\in M:\ f(m)=n\) می‌خوانیم.
        نماد دست‌نویس \(g(m)=n\) است؛ این اصلاح هنوز در انتظار تأیید است.
    \end{alertblock}
\end{frame}

% Source: sources/lectures/16 Nakayama Lemma and Exact sequences .pdf | PDF page: 5 | item: conclusion of the injectivity proof
% AMBIGUITY[L16-Q01]: proof continuation uses the explicitly provisional preimage f(m)=n.
\begin{frame}[t]{لم پنج‌تایی کوتاه: پایان اثبات یک‌به‌یکی}
    \label{lec16:five-injective-end}
    با خوانش موقتِ \(f(m)=n\) در گام قبل، جابجاییِ مربع چپ می‌دهد
    \[
        \varphi\alpha(m)=\beta f(m)=\beta(n)=0.
    \]
    از دقت سطر پایین، \(\varphi\) یک‌به‌یک است؛ پس
    \[
        \alpha(m)\in\ker\varphi=0\quad\Longrightarrow\quad\alpha(m)=0.
    \]
    چون \(\alpha\) یک‌به‌یک است،
    \[
        m=0\quad\Longrightarrow\quad n=f(m)=0.
    \]
    بنابراین \(\beta\) یک‌به‌یک است.
    \medskip

    \alert{این بازنویسیِ اثبات به خوانش موقتِ \lr{L16-Q01} وابسته است.}
\end{frame}

% Source: sources/lectures/16 Nakayama Lemma and Exact sequences .pdf | PDF page: 5 | item: surjectivity proof, choices of k and n
\begin{frame}[t]{لم پنج‌تایی کوتاه: پوشاییِ \(\beta\)}
    \label{lec16:five-surjective-choices}
    فرض کنید \(\alpha,\gamma\) پوشا باشند و \(n'\in N'\) دلخواه باشد.
    چون \(\gamma\) پوشاست،
    \[
        \psi(n')\in K'\quad\Longrightarrow\quad
        \exists k\in K:\ \gamma(k)=\psi(n').
    \]
    از دقت سطر بالا، \(g\) پوشاست؛ پس
    \[
        \exists n\in N:\quad g(n)=k.
    \]
    جابجاییِ مربع راست می‌دهد
    \[
        \psi\beta(n)=\gamma g(n)=\gamma(k)=\psi(n').
    \]
    در نتیجه
    \[
        n'-\beta(n)\in\ker\psi=\operatorname{Im}\varphi.
    \]
\end{frame}

% Source: sources/lectures/16 Nakayama Lemma and Exact sequences .pdf | PDF page: 6 | item: completion of the surjectivity proof
\begin{frame}[t]{لم پنج‌تایی کوتاه: پایان اثبات پوشایی}
    \label{lec16:five-surjective-end}
    عنصری \(m'\in M'\) وجود دارد که
    \[
        \varphi(m')=n'-\beta(n).
    \]
    چون \(\alpha\) پوشاست، برای یک \(m\in M\)، داریم \(\alpha(m)=m'\).
    بنابراین
    \[
        n'-\beta(n)=\varphi(m')=\varphi\alpha(m)=\beta f(m),
    \]
    و در نتیجه
    \[
        n'=\beta(n)+\beta f(m)=\beta\bigl(n+f(m)\bigr).
    \]
    پس \(\beta\) پوشاست. \hfill\(\qed\)
\end{frame}

% Source: sources/lectures/16 Nakayama Lemma and Exact sequences .pdf | PDF page: 6 | item: part iii follows from parts i and ii
% AMBIGUITY[L16-Q01]: the isomorphism conclusion cites part i, whose written proof uses the provisional map-name correction.
\begin{frame}[t]{حالت یکریختی در لم پنج‌تایی کوتاه}
    \label{lec16:five-isomorphism}
    \begin{theorem}
        در نمودار جابجاییِ دو رشتهٔ دقیق کوتاه، اگر \(\alpha\) و \(\gamma\)
        یکریختی باشند، آنگاه \(\beta\) نیز یکریختی است.
    \end{theorem}
    \textbf{اثبات جزوه.} از قسمت‌های اول و دوم، \(\beta\) هم یک‌به‌یک و هم پوشاست.
    \hfill\(\qed\)
    \medskip

    \begin{alertblock}{یادداشت بازنویسیِ اثبات: \lr{L16-Q01}}
        متن اثبات قسمت اول با خوانش موقتِ نام نگاشت ثبت شده است؛
        این ارجاع، تأیید نهاییِ آن نماد دست‌نویس نیست.
    \end{alertblock}
\end{frame}

% Source: sources/lectures/16 Nakayama Lemma and Exact sequences .pdf | PDF page: 6 | item: definition of equivalence of short exact sequences
\begin{frame}[t]{یکریختیِ دو رشتهٔ دقیق کوتاه}
    \label{lec16:equivalent-sequences}
    \begin{definition}
        دو رشتهٔ دقیق کوتاه را \textbf{یکریخت} می‌گوییم هرگاه یکریختی‌های
        \(\alpha,\beta,\gamma\) وجود داشته باشند که نمودار زیر جابجایی شود:
        \[
        \begin{tikzcd}[ampersand replacement=\&,column sep=large,row sep=large]
            0\arrow[r]\& M\arrow[r,"f"]\arrow[d,"\alpha"']
            \&N\arrow[r,"g"]\arrow[d,"\beta"]
            \&K\arrow[r]\arrow[d,"\gamma"]\&0\\
            0\arrow[r]\& M'\arrow[r,"\varphi"']
            \&N'\arrow[r,"\psi"']\&K'\arrow[r]\&0
        \end{tikzcd}
        \]
    \end{definition}
\end{frame}

% Source: sources/lectures/16 Nakayama Lemma and Exact sequences .pdf | PDF page: 7 | item: splitting lemma, three equivalent conditions
\begin{frame}[t]{قضیهٔ شکافتن رشتهٔ دقیق کوتاه}
    \label{lec16:splitting-lemma}
    رشتهٔ دقیق کوتاه زیر را در نظر بگیرید:
    \[
        0\longrightarrow M\xrightarrow{\varphi}K
        \xrightarrow{\psi}N\longrightarrow0.
    \]
    \begin{theorem}
        موارد زیر هم‌ارزند:
        \begin{enumerate}
            \item همریختیِ \(h:N\to K\) وجود دارد که \(\psi h=1_N\).
            \item همریختیِ \(g:K\to M\) وجود دارد که \(g\varphi=1_M\).
            \item رشتهٔ بالا با رشتهٔ دقیق کوتاهِ
            \[
                0\longrightarrow M\xrightarrow{i_1}M\oplus N
                \xrightarrow{\pi_2}N\longrightarrow0
            \]
            یکریخت است.
        \end{enumerate}
    \end{theorem}
\end{frame}

% Source: sources/lectures/16 Nakayama Lemma and Exact sequences .pdf | PDF page: 7 | item: proof i implies iii, definition of lambda and diagram
\begin{frame}[t]{از وارون راست تا شکافتن رشته}
    \label{lec16:split-right-map}
    فرض کنید \(\psi h=1_N\). همریختی زیر را تعریف می‌کنیم:
    \[
        \lambda:M\oplus N\longrightarrow K,
        \qquad \lambda(m,n)=\varphi(m)+h(n).
    \]
    نمودار مورد نظر چنین است:
    \[
    \begin{tikzcd}[ampersand replacement=\&,column sep=large,row sep=large]
        0\arrow[r]\& M\arrow[r,"i_1"]\arrow[d,"1_M"']
        \&M\oplus N\arrow[r,"\pi_2"]\arrow[d,"\lambda"]
        \&N\arrow[r]\arrow[d,"1_N"]\&0\\
        0\arrow[r]\& M\arrow[r,"\varphi"']
        \&K\arrow[r,"\psi"']\&N\arrow[r]\&0
    \end{tikzcd}
    \]
\end{frame}

% Source: sources/lectures/16 Nakayama Lemma and Exact sequences .pdf | PDF pages: 7-8 | item: commutative squares for lambda and application of short five lemma
% AMBIGUITY[L16-Q01]: source invokes the short five lemma; its proof-transcription caveat is retained in the dependency note.
\begin{frame}[t]{جابجاییِ نمودارِ \(\lambda\)}
    \label{lec16:split-right-squares}
    برای مربع چپ،
    \[
        \lambda i_1(m)=\lambda(m,0)=\varphi(m)=\varphi1_M(m).
    \]
    برای مربع راست،
    \begin{align*}
        \psi\lambda(m,n)
            &=\psi\bigl(\varphi(m)+h(n)\bigr)\\
            &=\psi\varphi(m)+\psi h(n)=0+n=n\\
            &=1_N\pi_2(m,n).
    \end{align*}
    پس نمودار جابجایی است. چون \(1_M,1_N\) یکریختی‌اند، از لم پنج‌تایی کوتاه،
    \(\lambda\) یکریختی است و شرط سوم برقرار می‌شود.
    \medskip

    {\footnotesize ارجاع به اثبات لم پیشین با یادداشت موقتِ \lr{L16-Q01} ثبت شده است.}
\end{frame}

% Source: sources/lectures/16 Nakayama Lemma and Exact sequences .pdf | PDF page: 8 | item: proof iii implies i, equivalence diagram
\begin{frame}[t]{از یکریختیِ رشته‌ها تا وارون راست}
    \label{lec16:split-equivalence-right}
    فرض کنید شرط سوم برقرار باشد؛ پس یکریختی‌های \(\theta,\eta,T\) نمودار زیر را جابجایی می‌کنند:
    \[
    \begin{tikzcd}[ampersand replacement=\&,column sep=large,row sep=large]
        0\arrow[r]\& M\arrow[r,"i_1"]\arrow[d,"\theta"']
        \&M\oplus N\arrow[r,"\pi_2"]\arrow[d,"\eta"]
        \&N\arrow[r]\arrow[d,"T"]\&0\\
        0\arrow[r]\& M\arrow[r,"\varphi"']
        \&K\arrow[r,"\psi"']\&N\arrow[r]\&0
    \end{tikzcd}
    \]
    از مربع راست،
    \[
        T\pi_2=\psi\eta
        \quad\Longrightarrow\quad
        T\pi_2i_2=\psi\eta i_2.
    \]
    چون \(\pi_2i_2=1_N\)، داریم \(T=\psi\eta i_2\).
\end{frame}

% Source: sources/lectures/16 Nakayama Lemma and Exact sequences .pdf | PDF page: 8 | item: construction of the right inverse of psi
\begin{frame}[t]{ساختن وارون راستِ \(\psi\)}
    \label{lec16:split-right-inverse}
    از رابطهٔ \(T=\psi\eta i_2\) و یکریختی بودن \(T\)،
    \[
        \psi\eta i_2T^{-1}=1_N.
    \]
    در نتیجه با انتخاب
    \[
        h=\eta i_2T^{-1}:N\longrightarrow K,
    \]
    به دست می‌آوریم
    \[
        \psi h=1_N.
    \]
    پس شرط اول برقرار است. \hfill\(\qed\)
\end{frame}

% Source: sources/lectures/16 Nakayama Lemma and Exact sequences .pdf | PDF page: 8 | item: proof iii implies ii, reverse equivalence diagram and construction of left inverse
\begin{frame}[t]{از یکریختیِ رشته‌ها تا وارون چپ}
    \label{lec16:split-left-inverse}
    با برعکس کردن یکریختی‌های رشته‌ها، نمودار زیر جابجایی است:
    \[
    \begin{tikzcd}[ampersand replacement=\&,column sep=large,row sep=large]
        0\arrow[r]\& M\arrow[r,"\varphi"]\arrow[d,"\eta_1"']
        \&K\arrow[r,"\psi"]\arrow[d,"\eta_2"]
        \&N\arrow[r]\arrow[d,"\eta_3"]\&0\\
        0\arrow[r]\& M\arrow[r,"i_1"']
        \&M\oplus N\arrow[r,"\pi_2"']\&N\arrow[r]\&0
    \end{tikzcd}
    \]
    از مربع چپ،
    \[
        \eta_2\varphi=i_1\eta_1
        \quad\Longrightarrow\quad
        \pi_1\eta_2\varphi=\eta_1.
    \]
    بنابراین
    \[
        (\eta_1^{-1}\pi_1\eta_2)\varphi=1_M.
    \]
    پس \(g=\eta_1^{-1}\pi_1\eta_2\) وارون چپِ \(\varphi\) است.
\end{frame}

% Source: sources/lectures/16 Nakayama Lemma and Exact sequences .pdf | PDF page: 9 | item: proof ii implies iii, lambda-prime and diagram; commutativity referred to the preceding case
% AMBIGUITY[L16-Q01]: the final isomorphism step invokes the short five lemma, with its provisional proof-transcription caveat.
\begin{frame}[t]{از وارون چپ تا شکافتن رشته}
    \label{lec16:split-left-map}
    فرض کنید \(g:K\to M\) و \(g\varphi=1_M\). مانند حالت قبل تعریف می‌کنیم:
    \[
        \lambda':K\longrightarrow M\oplus N,
        \qquad k\longmapsto\bigl(g(k),\psi(k)\bigr).
    \]
    همانند حالت قبل، نمودار زیر جابجایی است:
    \[
    \begin{tikzcd}[ampersand replacement=\&,column sep=large,row sep=large]
        0\arrow[r]\& M\arrow[r,"\varphi"]\arrow[d,"1_M"']
        \&K\arrow[r,"\psi"]\arrow[d,"\lambda'"]
        \&N\arrow[r]\arrow[d,"1_N"]\&0\\
        0\arrow[r]\& M\arrow[r,"i_1"']
        \&M\oplus N\arrow[r,"\pi_2"']\&N\arrow[r]\&0
    \end{tikzcd}
    \]
    از لم پنج‌تایی کوتاه، \(\lambda'\) یکریختی است؛ پس شرط سوم برقرار است.
    {\footnotesize یادداشتِ بازنویسی اثبات لم پیشین: \lr{L16-Q01}.}
\end{frame}

% Source: sources/lectures/16 Nakayama Lemma and Exact sequences .pdf | PDF page: 9 | item: split exact sequence definition and direct sum conclusion
\begin{frame}[t]{رشتهٔ دقیق شکافته‌شده}
    \label{lec16:split-exact-definition}
    \begin{definition}
        رشتهٔ دقیق کوتاهِ
        \[
            0\longrightarrow M\xrightarrow{\varphi}K
            \xrightarrow{\psi}N\longrightarrow0
        \]
        را \textbf{شکافته‌شده} می‌گوییم هرگاه یکی از شرایط قضیهٔ قبل در آن صدق کند.
    \end{definition}
    در این صورت
    \[
        K\simeq M\oplus N.
    \]
\end{frame}
